\section{Acceleration}

\subsection{Why velocity is not the end of the description}

Velocity describes how position changes. But velocity itself may also change. A body may move faster and faster, slow down, or move along a curve while its direction of motion changes.

Acceleration describes the change of velocity with respect to time.

\subsection{Average acceleration}

Suppose the velocity at time \(t_1\) is \(\mathbf{v}(t_1)\), and the velocity at time \(t_2\) is \(\mathbf{v}(t_2)\), with \(t_1\neq t_2\). The change in velocity is
\[
\Delta\mathbf{v}=\mathbf{v}(t_2)-\mathbf{v}(t_1).
\]

\begin{definitionbox}
The average acceleration of a moving point on the interval \([t_1,t_2]\), with \(t_1\neq t_2\), is
\[
\mathbf{a}_{\mathrm{avg}}
=
\frac{\Delta\mathbf{v}}{\Delta t}
=
\frac{\mathbf{v}(t_2)-\mathbf{v}(t_1)}{t_2-t_1}.
\]
\end{definitionbox}

Average acceleration describes how velocity changes over a time interval.

\subsection{Instantaneous acceleration}

\begin{definitionbox}
Let \(\mathbf{v}\) be differentiable at time \(t\). The instantaneous acceleration at time \(t\) is
\[
\mathbf{a}(t)=\frac{d\mathbf{v}}{dt}(t).
\]
\end{definitionbox}

Since velocity is the derivative of position,
\[
\mathbf{v}(t)=\frac{d\mathbf{r}}{dt},
\]
we also have
\[
\mathbf{a}(t)=\frac{d^2\mathbf{r}}{dt^2}.
\]
Thus acceleration is the second derivative of position with respect to time.

In coordinates,
\[
\mathbf{a}(t)=
\left(
\frac{d^2x}{dt^2},
\frac{d^2y}{dt^2},
\frac{d^2z}{dt^2}
\right).
\]
In two dimensions,
\[
\mathbf{a}(t)=
\left(
\frac{d^2x}{dt^2},
\frac{d^2y}{dt^2}
\right),
\]
and in one dimension,
\[
a(t)=\frac{d^2x}{dt^2}.
\]

\subsection{Acceleration is also a vector}

Acceleration is a vector. It has magnitude and direction. It describes how the velocity vector changes.

Since velocity is a vector, it can change in magnitude, in direction, or in both. Therefore a body may be accelerating even if its speed is constant, provided that the direction of velocity changes.

\begin{remarkbox}
Constant speed does not necessarily mean zero acceleration. If the direction of velocity changes, the velocity vector changes, and the acceleration is not zero.
\end{remarkbox}

\subsection{Examples}

\begin{examplebox}
Consider one-dimensional motion
\[
x(t)=x_0+at,
\]
where \(x_0\) and \(a\) are constants. Then
\[
v(t)=a.
\]
The acceleration is
\[
a_{\mathrm{kin}}(t)=\frac{dv}{dt}=0.
\]
Here \(a_{\mathrm{kin}}\) denotes acceleration, while \(a\) is a parameter in the formula for \(x(t)\).
\end{examplebox}

\begin{examplebox}
For
\[
\mathbf{r}(t)=(t,t^2),
\]
we found
\[
\mathbf{v}(t)=(1,2t).
\]
Therefore
\[
\mathbf{a}(t)=\frac{d\mathbf{v}}{dt}=(0,2).
\]
The acceleration is constant and points in the positive \(y\)-direction.
\end{examplebox}

\begin{examplebox}
For circular motion
\[
\mathbf{r}(t)=(R\cos t,R\sin t),
\]
we found
\[
\mathbf{v}(t)=(-R\sin t,R\cos t).
\]
Differentiating again gives
\[
\mathbf{a}(t)=(-R\cos t,-R\sin t).
\]
Since
\[
\mathbf{r}(t)=(R\cos t,R\sin t),
\]
we get
\[
\mathbf{a}(t)=-\mathbf{r}(t).
\]
The acceleration points toward the center of the circle.
\end{examplebox}

This is a major conceptual point. In uniform circular motion, speed may be constant, but acceleration is not zero, because the direction of velocity is changing.

\subsection{Units of acceleration}

Acceleration has units of velocity divided by time. If velocity is measured in metres per second and time in seconds, acceleration is measured in
\[
\mathrm{m/s^2}.
\]
This also follows from
\[
\mathbf{a}(t)=\frac{d^2\mathbf{r}}{dt^2}.
\]
Position contributes metres, and two differentiations with respect to time contribute two factors of seconds in the denominator. The unit \(\mathrm{m/s^2}\) means metres per second per second.
